% One turn of the loop in Algorithm 2, shown on the graph rather than on the
% drawing: the region arrives with every crossing edge present and leaves with
% exactly one of them kept.
\begin{figure}[htbp]
  \centering
  \begin{subfigure}[b]{0.47\textwidth}\centering
    \begin{tikzpicture}[scale=0.82]
      \foreach \x in {0,1,2,3,4} \foreach \y in {0,1,2,3}
        \node[node vertex] (u\x\y) at (\x,-\y) {};
      \foreach \x/\xn in {0/1,1/2,2/3,3/4} \foreach \y in {0,1,2,3}
        \draw[black!45] (u\x\y) -- (u\xn\y);
      \foreach \y/\yn in {0/1,1/2,2/3} \foreach \x in {0,1,2,3,4}
        \draw[black!45] (u\x\y) -- (u\x\yn);
    \end{tikzpicture}
    \caption{On entry the region is the full grid graph on
    $\region{a_0,a_1}{b_0,b_1}$: connected, and full of cycles.}
    \label{fig:partition-before}
  \end{subfigure}\hfill
  \begin{subfigure}[b]{0.47\textwidth}\centering
    \begin{tikzpicture}[scale=0.82]
      \foreach \x in {0,1,2,3,4} \foreach \y in {0,1,2,3}
        \node[node vertex] (w\x\y) at (\x,-\y) {};
      \foreach \x/\xn in {0/1,1/2,3/4} \foreach \y in {0,1,2,3}
        \draw[black!45] (w\x\y) -- (w\xn\y);
      \foreach \y/\yn in {0/1,1/2,2/3} \foreach \x in {0,1,2,3,4}
        \draw[black!45] (w\x\y) -- (w\x\yn);
      \draw[door] (w21) -- (w31);
      \draw[wall] (2.5,0.45) -- (2.5,-0.75);
      \draw[wall] (2.5,-1.25) -- (2.5,-3.45);
      \node[font=\tiny, wallred, above] at (2.5,0.45) {wall at $j=2$};
      \node[font=\tiny, doorgreen, left] at (-0.2,-1) {door at $i{=}1$};
    \end{tikzpicture}
    \caption{On exit every edge crossing column boundary $j$ is gone except the
    one at row $i$.}
    \label{fig:partition-after}
  \end{subfigure}
  \caption{One partition step, with $t=1$ (a vertical wall). The loop walks
  $k$ over the rows $b_0..b_1$ of the region and deletes the edge
  $(j,k)\!-\!(j{+}1,k)$ for every $k\neq i$. The two halves stay connected
  internally, and the single surviving edge is the only way between them ---
  which is exactly the induction step of
  Proposition~\ref{prop:generation-is-a-tree}.}
  \label{fig:one-partition-step}
\end{figure}
